\section{System Architecture}
\label{sec:architecture}

This chapter describes how the complete system is organised: the three
physical tiers on which computation runs, the data flow within and across
those tiers, the inventory of software components, and the design rules
that shaped every implementation decision. The architecture is deliberately
asymmetric. Almost all intelligence lives on one small computer that is
physically co-located with the rover fleet, while the rovers themselves
stay minimal and the cloud is used only where its latency does not matter.

\subsection{The Three-Tier Deployment Model}
\label{sec:architecture-tiers}

We adopt the mist/fog/cloud hierarchy that is standard in edge computing
for robotics, and map our pipeline onto it as follows.

\paragraph{Mist tier: the rovers.}
Each rover carries an ESP-NOW radio and a microcontroller with limited
memory and no floating-point acceleration to speak of. The rover's
responsibilities are confined to sensing (link quality to neighbours,
battery state, odometry), executing the motion commands implied by an
assigned task, and reporting task outcomes back to the fog node. No
optimisation, no graph reasoning, and no learning happens on a rover.
This choice reflects the hardware we target: a rover that must be cheap
enough to lose in the field cannot host an annealer or a neural network.
ESP-NOW provides short-range, low-power, peer-to-peer frames without the
overhead of an association-based protocol, which suits the dense, rapidly
changing local topologies we study.

\paragraph{Fog tier: the fleet brain.}
A single Raspberry Pi class computer sits within radio range of the fleet
and runs the entire per-cycle decision loop: it assembles the candidate
graph from rover reports, formulates and solves the constrained QUBO edge
pruning problem, scores rovers with the trust-aware GNN on the pruned
graph, and performs the greedy capacity-aware task assignment. This tier
is the reason for the most consequential engineering constraint of the
project: \emph{all per-cycle numerical code is pure NumPy over tiny
matrices}. A cycle must complete in well under the control period on a
few-watt ARM processor, so there is no room for a deep learning
framework, a sparse linear algebra stack, or a commercial solver. Every
matrix in the inference path is small enough that dense NumPy with a
handful of temporaries is both the simplest and the fastest honest
option. The message-passing round of the GNN is a single adjacency
product followed by an elementwise nonlinearity, and the projected edge
annealer manipulates flat bit vectors and small dictionaries rather than
any external optimisation library.

\paragraph{Cloud tier: training, amortised solves, analytics.}
Anything whose latency does not couple into the control loop is pushed to
the cloud tier. Three workloads live there. First, offline training of the
GNN readout layer on logged outcomes; we report in Chapter 6 that this
training did not beat the fixed readout on held-out data, but the
infrastructure remains part of the architecture. Second, amortised heavy
solves: the pruning QUBO can be re-solved every $k$ cycles rather than
every cycle, and heavy verification solves (the exact MILP baseline, or
cross-checks of annealer output) are cloud jobs, not field jobs. Third,
fleet-level analytics over mission logs.

Figure~\ref{fig:tiers} summarises the deployment and the data that crosses
tier boundaries.

\begin{figure}[ht]
\centering
\resizebox{0.95\textwidth}{!}{%
\begin{tikzpicture}[
  font=\small,
  tier/.style={draw, rounded corners=2pt, thick, inner sep=8pt},
  title/.style={font=\bfseries\small},
  comp/.style={draw, fill=black!4, rounded corners=1pt,
               inner sep=3pt, font=\footnotesize, align=center},
  flow/.style={-{Stealth[length=2.5mm]}, thick},
  lbl/.style={font=\scriptsize, align=center, inner sep=2pt},
]
% --- Mist tier box ---
\node[comp] (rovers) at (0,1.0) {rover 1 \ldots\ rover $n$\\[-1pt]
  \scriptsize ESP-NOW radio, sensors, actuators};
\node[comp, below=4pt of rovers] (mistnote) {sense links, battery, odometry\\
  execute tasks, report outcomes};
\begin{scope}[on background layer]
  \node[tier, fit=(rovers)(mistnote), label={[title]north:Mist tier (rovers)}] (mist) {};
\end{scope}

% --- Fog tier box ---
\node[comp, right=3.4cm of rovers] (qubo) {QUBO builder\\+ projected edge annealer};
\node[comp, below=4pt of qubo] (gnn) {trust model + trust-aware GNN};
\node[comp, below=4pt of gnn] (assign) {greedy capacity-aware assigner};
\begin{scope}[on background layer]
  \node[tier, fit=(qubo)(gnn)(assign),
        label={[title]north:Fog tier (Raspberry Pi class)}] (fog) {};
\end{scope}

% --- Cloud tier box ---
\node[comp, right=3.4cm of qubo] (train) {readout training\\(offline)};
\node[comp, below=4pt of train] (heavy) {amortised heavy solves\\exact MILP verification};
\node[comp, below=4pt of heavy] (analytics) {fleet-level analytics};
\begin{scope}[on background layer]
  \node[tier, fit=(train)(heavy)(analytics),
        label={[title]north:Cloud tier (offline)}] (cloud) {};
\end{scope}

% --- Flows: mist <-> fog ---
\draw[flow] ($(mist.east)+(0,0.35)$) -- node[lbl, above]{state reports\\(links, battery, outcomes)}
  ($(fog.west)+(0,0.35)$);
\draw[flow] ($(fog.west)-(0,0.35)$) -- node[lbl, below]{task assignments}
  ($(mist.east)-(0,0.35)$);

% --- Flows: fog <-> cloud ---
\draw[flow] ($(fog.east)+(0,0.5)$) -- node[lbl, above]{mission logs, QUBO\\instances, outcomes}
  ($(cloud.west)+(0,0.5)$);
\draw[flow] ($(cloud.west)-(0,0.5)$) -- node[lbl, below]{readout weights\\(rare updates)}
  ($(fog.east)-(0,0.5)$);
\end{tikzpicture}}
\caption{Three-tier deployment. The mist tier (rovers) senses and acts;
the fog tier runs the complete per-cycle decision loop on a Raspberry Pi
class computer using pure NumPy; the cloud tier performs offline training,
amortised or verification-grade solves, and fleet analytics.}
\label{fig:tiers}
\end{figure}

\subsection{Per-Cycle Data Flow}
\label{sec:architecture-flow}

One control cycle proceeds through six stages, which we describe in order,
noting what crosses a tier boundary at each step.

\paragraph{Stage 1: graph assembly (mist to fog).}
Each rover reports its current estimate of link quality to every
neighbour it can hear, its battery level, its compute availability, and
its position estimate. These reports travel over ESP-NOW to the fog node,
which merges them with its own per-node records (trust, freshness of last
contact, task criticality, processing cost) to build the full dynamic
candidate graph: every admissible rover pair is an edge, and every edge
carries a feature vector in $[0,1]^7$ covering trust, link quality,
battery, proximity, freshness, task criticality, and processing cost.

\paragraph{Stage 2: constrained QUBO pruning (fog).}
The candidate graph is denser than the GNN needs, and dense graphs dilute
the trust signal and increase inference cost. The QUBO builder encodes
edge selection as a binary quadratic program: a linear utility per edge
(the learned-weight combination of the seven features), soft penalties
for graph size and for churn against the previous cycle's decisions, and
hard constraints with dominant penalty weights that enforce per-node
degree limits, a global processing budget, and single-commodity flow
connectivity from a root to every required node. The production solver,
the projected edge annealer, moves directly in the space of feasible edge
sets: it toggles or swaps edges and repairs each candidate via an exact
slack-and-flow recomputation, so infeasible states are never committed.
The output is a pruned subgraph plus, for auditability, the raw solution
vector and a list of any violated constraints.

\paragraph{Stage 3: trust-aware scoring (fog).}
The trust model maintains an exponentially weighted moving average of
per-node reliability, with ageing so that stale evidence decays. The GNN
takes the pruned adjacency, per-node features (trust, battery, available
compute, capability), and performs one message-passing round with a
sigmoid nonlinearity, then a linear readout scores each rover against the
current task requirements. All of this is dense NumPy on matrices of
order $n \times n$ with $n$ the fleet size, comfortably inside the fog
compute budget.

\paragraph{Stage 4: assignment (fog to mist).}
A greedy capacity-aware matcher assigns tasks to the highest-scoring
rovers that still have free capacity. The assignments are radioed down to
the rovers, which execute them.

\paragraph{Stage 5: outcome sampling (mist).}
Task outcomes are realised on the rovers. In simulation, outcomes are
sampled from a hidden, spatially correlated reliability field, and the
probability of successful delivery depends causally on the pruned
topology through the paths the assignment actually uses; this design,
introduced after reviewer pressure, prevents the evaluation from being
merely correlational.

\paragraph{Stage 6: trust update (fog, then feedback).}
Reported outcomes update the trust model, closing the loop: the updated
trust values are inputs to the edge features of the next cycle. Mission
logs, QUBO instances, and outcome traces are streamed or batched to the
cloud tier for the offline workloads described above; the only
cloud-to-fog traffic in normal operation is an occasional readout-weight
update.
\subsection{Component Inventory}
\label{sec:architecture-components}

Table~\ref{tab:components} lists the components of the system by role.
Two deliberate asymmetries are worth flagging. First, two annealers exist:
a bit-level simulated annealer that manipulates raw bit strings and is
used \emph{only} as a benchmark subject, and the projected edge annealer
that is the production solver. The benchmark exists precisely because the
bit-level solver performs poorly (it cannot cross penalty barriers and
returned infeasible states on most benchmark instances), which is itself
one of our audit findings. Second, the exact MILP baseline (CBC via PuLP)
is faster and exact on our instance sizes, which we report honestly as a
negative result for QUBO on classical hardware; it serves as the ground
truth oracle for the solver audit.

\begin{table}[ht]
\centering
\caption{Component inventory. ``Tier'' indicates where the component runs
in deployment; benchmark and verification components are cloud-side or
development-side tools.}
\label{tab:components}
{\small
\begin{tabular}{@{}llll@{}}
\toprule
Component & Role & Tier & Technology \\
\midrule
QUBO builder & constrained edge-pruning encoding & fog & NumPy \\
Bit-level annealer & benchmark subject only & cloud/dev & NumPy \\
Projected edge annealer & production QUBO solver & fog & NumPy \\
Trust model & EWMA reliability with ageing & fog & NumPy \\
Trust-aware GNN & rover scoring on pruned graph & fog & NumPy \\
Greedy assigner & capacity-aware task matching & fog & NumPy \\
Exact MILP baseline & ground-truth oracle, audit & cloud & CBC via PuLP \\
Scaling runner & fleet-size experiments & cloud & NumPy \\
Hostile test suite & 491 checks, sections A to N & dev & Python stdlib \\
Causal simulation harness & end-to-end evaluation, statistics & cloud & NumPy \\
Physics integration & realistic rover dynamics & cloud & Gazebo, ROS 2, Docker \\
Fleet console & live visualisation and tuning & dev & three.js, React, Python \\
\bottomrule
\end{tabular}
}
\end{table}

The \textbf{causal simulation harness} drives the whole pipeline
end-to-end against the hidden reliability field and hosts the pruning
baselines (no pruning, threshold pruning, top-$k$ pruning, full graph)
and assignment baselines (trust-only, no-graph, random) used in the
ablation studies. It reports paired bootstrap confidence intervals and
permutation $p$-values so that every comparative claim carries an
uncertainty statement.

The \textbf{physics integration} replaces the abstract kinematics of the
simulator with Gazebo Harmonic under ROS 2 Jazzy, packaged in Docker so
that the planner node runs unmodified against simulated physics; this
validates that the fog-side loop tolerates realistic sensing noise and
actuation latency.

The \textbf{fleet console} is a browser application: a custom three.js
rendering engine (instanced meshes, a single instanced shader draw for
link arcs) with a React interface and Zustand state management, served by
a Python standard-library HTTP backend that steps live simulation
sessions and exposes in-place parameter tuning. It exists for
demonstration and for interactive exploration of solver parameters, and
its comparison menu names the baselines that beat our pipeline rather
than hiding them.

The \textbf{hostile test suite} deserves special mention as an
architectural component rather than an afterthought. Its 491 checks,
organised into sections A to N, include an independent re-derivation of
the QUBO energy that is cross-verified against the quadratic form itself,
brute-force enumeration of small instances as ground truth, exact MILP
cross-checks on larger instances, and sensitivity sweeps over annealer
hyperparameters and penalty scales.

\subsection{Design Rules}
\label{sec:architecture-rules}

Four rules, adopted early and enforced throughout, explain most of the
concrete choices above.

\paragraph{Rule 1: the fog compute budget is real.}
Every per-cycle component must run in milliseconds on a Raspberry Pi
class machine. This rules out deep learning frameworks (hence a GNN that
is one dense message-passing round in NumPy), rules out heavy sparse or
optimisation libraries in the inference path (hence a hand-rolled
projected annealer), and motivates QUBO pruning in the first place: the
pruned graph is cheaper to score, so the pruner partially pays for
itself.

\paragraph{Rule 2: no heavy dependencies anywhere.}
The numerical stack depends on NumPy alone. The exact baseline uses PuLP
with CBC, both permissively licensed and installable without system
libraries; notably, we avoid SciPy entirely. The web backend uses only
the Python standard library. The physics simulator is containerised with
Docker precisely so that its heavier ROS 2 dependency chain is isolated
from the numerical core.

\paragraph{Rule 3: every claim is paired with a test.}
The 491-check hostile suite is co-developed with the code, not appended
to it. The solver exactness rates quoted in later chapters are enforced
by floor checks in the suite, the simulation statistics are regression
checked, and new components are expected to arrive with their own hostile
tests.

\paragraph{Rule 4: negative results are kept.}
The architecture is an audit, not only a proposal, and the system is
built to surface unfavourable comparisons rather than suppress them. The
exact MILP baseline that dominates the QUBO solver on classical hardware
is a first-class component; the readout-training null result is reported
with its $p$-value; the pruner comparison in which threshold pruning wins
on raw mission success at the ten-rover scale appears in the results with
its significance test. We regard these findings as outputs of the system,
on equal footing with its successes.

\subsection{Summary}
The system places a complete decision loop on one modest fog computer by
discipline rather than by excess compute: a minimal mist tier, a
NumPy-only fog tier running pruning, scoring, and assignment, and a
cloud tier reserved for training, verification-grade solves, and
analytics. The
surrounding tooling, from the MILP oracle to the hostile suite to the
fleet console, exists to make the central claims falsifiable and the
negative results impossible to misplace.
